\snsection{ASPECTOS FINANCIEROS}
\justifying

La gestión eficiente de los recursos financieros es un aspecto crítico en el desarrollo de proyectos, especialmente cuando se involucran componentes especializados y procesos de fabricación personalizados. La \cref{tab:aspectos_financieros:gastos} presenta un desglose detallado de los gastos del proyecto, convertidos al valor de venta del dólar del Banco Nación en cada fecha. Esta información no solo permite un seguimiento preciso de los costos, sino que también sirve como base para evaluar la viabilidad económica del proyecto y optimizar futuras asignaciones presupuestarias.

% Mateo: hay que ver de qué forma queda más linda la tabla, no sé si así está bien (ver descripción de cada gasto).
\begin{table}[H]
\centering
\caption{Gastos del proyecto}
\begin{tabularx}{\linewidth}{lXr}
\hline
\textbf{Fecha} & \textbf{Descripción}                                   & \textbf{Costo} \\ \hline
08/04/2025     & Componentes electrónicos varios en Electrocomponentes  & 19,16 USD      \\
23/06/2025     & Componentes electrónicos varios en it\&t               & 90,51 USD      \\
02/07/2025     & Placa de desarrollo ESP32-S3 N16R8                     & 6,88 USD       \\
16/07/2025     & Pantalla TFT LCD 2,9"                                  & 20,88 USD      \\
17/07/2025     & Placas adaptador M2 a PCIe X4                          & 6,43 USD       \\
22/07/2025     & Componentes electrónicos varios en Elemon              & 23,53 USD      \\
13/08/2025     & Transformador 220 V/48 V 50 VA                         & 41,51 USD      \\
19/08/2025     & Impresión 3D \#1                                       & 3,82 USD       \\
N/A            & Insumos varios                                         & 15,38 USD      \\ \hline
\textbf{Total} &                                                        & 228,1 USD     \\ \hline

\label{tab:aspectos_financieros:gastos}
\end{tabularx}
\end{table}

\begin{comment}
Describa aquí los costos de materiales del prototipo y recursos con los que cuenta.
Debe estar en la misma unidad monetaria que en aspectos de mercado.

No poner precio al proyecto, solamente costo de materiales.
\end{comment}